\section{Example 1: Simple Arithmetic}

\textbf{TypeScript equivalent:}
\begin{lstlisting}[style=ts]
function add(a: number, b: number): number {
  return a + b;
}
\end{lstlisting}

\textbf{\acr{LLVM} \acr{IR} generation:}
\begin{lstlisting}[style=ts]
import {
  LLVMCodegen,
  LLVMFunction,
  Parameter,
  BasicBlock,
  BinaryInst,
  RetInst
} from "@euriklis/llvm-ir";

const codegen = new LLVMCodegen("arithmetic");

const addFunc = new LLVMFunction({
  name: "add",
  returnType: LLVMCodegen.types.i32
});

addFunc.addParameter(new Parameter({
  type: LLVMCodegen.types.i32,
  name: "a"
}));
addFunc.addParameter(new Parameter({
  type: LLVMCodegen.types.i32,
  name: "b"
}));

const entry = new BasicBlock("entry");
entry.addInstruction(new BinaryInst({
  name: "result",
  opcode: LLVMCodegen.opcodes.add,
  type: LLVMCodegen.types.i32,
  lhs: "a",
  rhs: "b"
}));
entry.setTerminator(new RetInst({
  type: LLVMCodegen.types.i32,
  value: "result"
}));

addFunc.addBasicBlock(entry);
codegen.getModule().addFunction(addFunc);

console.log(codegen.toString());
\end{lstlisting}

\textbf{Generated \acr{LLVM} \acr{IR}:}
\begin{lstlisting}[style=ir]
target triple = "x86_64-unknown-linux-gnu"

define i32 @add(i32 %a, i32 %b) {
entry:
  %result = add i32 %a, %b
  ret i32 %result
}
\end{lstlisting}

\textbf{Key concepts:}
\begin{itemize}[noitemsep]
  \item \texttt{LLVMFunction}: Defines function signature
  \item \texttt{Parameter}: Typed function argument (automatically prefixed with \texttt{\%})
  \item \texttt{BasicBlock}: Container for instructions (must end with terminator)
  \item \texttt{BinaryInst}: Arithmetic operation (add, sub, mul, etc.)
  \item \texttt{RetInst}: Return value (terminates basic block)
\end{itemize}

